% chercher des documents LaTeX dans styles, corps et bib
\makeatletter\def\input@path{{styles/}{corps/}}\makeatother

% Utiliser le style rapport.cls
\documentclass{rapport}
\usepackage{fancyhdr}

\fancyhead[L,C]{}
\fancyhead[R]{Rapport de 1�re phase}
\fancyfoot[C]{Alexandre Froehlich \& Guillaume Leinen}
\fancyfoot[L]{U.E. 2.4}
\fancyfoot[R]{\thepage}
\renewcommand{\headrulewidth}{0.4pt}
\renewcommand{\footrulewidth}{0.4pt}


\title{Rapport de 1�re phase - UE 2.4.}

\author{Alexandre \textsc{Froehlich} \& Guillaume \textsc{Leinen}\\
  \texttt{alexandre.froehlich@ensta-bretagne.org}
  \texttt{guillaume.leinen@ensta-bretagne.org}}
\date{\today}
\doctype{Projet Info.\\ U.E. 2.4.}
\promo{FISE 2022}
\etablissement{\textsc{Ensta} Bretagne\\2, rue Fran�ois Verny\\
  29806 \textsc{Brest} cedex\\\textsc{France}\\Tel +33 (0)2 98 34 88 00\\ \url{www.ensta-bretagne.fr}}
\logoEcole{\includegraphics[height=4.2cm]{logo_ENSTA_Bretagne_Vertical_CMJN}}

% M�tadonn�es � include dans le pdf
\hypersetup{%
  pdftitle={Rapport premi�re phase},
  pdfauthor={Alexandre Froehlich, Guillaume Leinen},
  pdfkeywords={},
  bookmarksnumbered,
  pdfstartview={FitH},
  citecolor=blue,
  breaklinks=true
}%

% creer le glossaire
\makeglossaries
% creer l'index
\makeindex

\pagestyle{fancy}

\begin{document}
% D�but du pr�ambule
\frontmatter
% inclure la liste des acronymes
\newglossaryentry{iot}{type=glo,
  text={\textsc{IoT}},
  name={Internet of Things},
  description={Ensembles des objets dits "connect�s". Cela va du smartphone, � la t�l�vision, au frigo...}}

\newglossaryentry{ia}{type=glo,
	text={\textsc{IA}},
	name={Intelligence Artificielle},
	description={Algorithme ou machine qui a pour but de simuler une intelligence}}

\newglossaryentry{enigma}{type=glo,
	text={\textsc{Enigma}},
	name={Enigma},
	description={Machine servant au chiffrement et d�chiffrement de messages, d'informations principalement employ�e par les Allemands}}

\newglossaryentry{alphago}{type=glo,
	text={\textsc{AlphaGo}},
	name={AlphaGo},
	description={Programme informatique qui a pour but de jouer au Go, il a remport� de nombreux prix face � de grands joueurs humains}}

\newglossaryentry{resneurone}{type=glo,
	text={\textsc{R�seau Neuronal}},
	name={R�seau Neuronal},
	description={Programme ou syst�me que l'on va entra�ner � reconnaitre des motifs, des tendances, ...}}

\newglossaryentry{dataset}{type=glo,
	text={\textsc{dataset}},
	name={Dataset},
	description={Ensemble de donn�es. Peut �tre des images, des vid�os, du texte, ...}}

\newglossaryentry{pixel}{type=glo,
	text={\textsc{pixel}},
	name={Pixel},
	description={Unit� �l�mentaire composant une image, poss�de une unique couleur}}

\newglossaryentry{MNIST}{type=glo,
	text={\textsc{MNIST}},
	name={MNIST Dataset},
	description={Ensemble de donn�es du Mixed National Institute of Standards and Technology, un dataset contenant des chiffres �crits � la main}}

\newglossaryentry{CNN}{type=glo,
	text={\textsc{CNN}},
	name={Convolutional Neural Network},
	description={R�seau neuronal dans lequel le motif de connexion entre les neurones est inspir� du cortex visuel des animaux}}

\newglossaryentry{RCNN}{type=glo,
	text={\textsc{R-CNN}},
	name={Regions Convolutional Neural Network},
	description={Un CNN qui ne travaille que sur une r�gion donn�e d'une information en entr�e}}

\newglossaryentry{TPU}{type=glo,
	text={\textsc{TPU}},
	name={Tensor Processing Unit},
	description={Mat�riel sp�cifiquement con�u par Google pour calculer les tenseurs, structure math�matique au c?ur de ses algorithmes}}

\newglossaryentry{framework}{type=glo,
	text={\textsc{framework}},
	name={Framework},
	description={Cadre de travail, c'est souvent un ensemble de m�thodes et de fonctions d�finies dans une biblioth�que}}

\newglossaryentry{npl}{type=glo,
	text={\textsc{Natural Processing Language}},
	name={Natural Processing Language},
	description={Ensemble de techniques qui ont pour but d'harmoniser le langage humain et le langage machine}}

\newglossaryentry{appsuper}{type=glo,
	text={\textsc{apprentissage supervis�}},
	name={Apprentissage supervis�},
	description={Apprentissage de pr�diction avec un dataset d'entr�e et le r�sultat suppos� en sortie}}

\newglossaryentry{appnonsuper}{type=glo,
	text={\textsc{apprentissage non-supervis�}},
	name={Apprentissage non-supervis�},
	description={Apprentissage sans pr�diction, on laisse le r�seau neuronal "explorer"}}

\newglossaryentry{apprenfor}{type=glo,
	text={\textsc{apprentissage par renforcement}},
	name={Apprentissage par renforcement},
	description={Apprentissage par s�lection ou r�compense. On r�compense le r�seau neuronal pour les r�ponses les plus adapt�es qu'il donne}}

\newglossaryentry{deep}{type=glo,
	text={\textsc{deep learning}},
	name={Deep Learning},
	description={Apprentissage profond. Devenu un terme g�n�rique pour parler d'un algorithme poss�dant un r�seau neuronal}}

\newglossaryentry{transfer}{type=glo,
	text={\textsc{transfer learning}},
	name={Transfer Learning},
	description={Concept qui prend un r�seau neuronal pr�-entra�n� et qui l'exploite sur une application similaire sans passer par la phase d'apprentissage exclusif}}

%%% Local Variables: 
%%% mode: latex
%%% TeX-master: "../guide"
%%% End: 

% creer le titre ici
\maketitle

\clearpage

% table des matieres
\tableofcontents

\clearpage

% Partie principale
\mainmatter
% chargement du fichier corps.tex

\section{Description g�n�rale du probl�me}
\subsection{Introduction}

\subsection{Les diff�rentes pistes envisag�es}

\subsection{Diff�rentes hypoth�ses r�ductrices choisies}


\section{Application proprement dite}
\subsection{Pr�sentation g�n�rale du programme}
% Ins�rer un sch�ma de fonctionnement

\subsection{Description des principales classes/m�thodes}
% d�crire uniquement les m�canismes fondamentaux


\section{Conclusion}
\subsection{Diff�rents tests effectu�s}

\subsection{Limitations}

\subsection{Apports personnels par rapport au sujet initial}

\subsection{Perspectives d'am�lioration}

\section{Introduction}
test



\clearpage 
\appendix
\pagestyle{empty}
% \section{A ajouter s'il le faut}


%\clearpage
%\phantomsection
%\label{sec:index}
%\addcontentsline{toc}{section}{Index}
%\printindex

%\printglossaries


%\clearpage

\end{document}

%%% Local Variables: 
%%% mode: latex
%%% TeX-master: t
%%% End: 
